\subsection{Image operators}\label{sec:operators}

The operator chain holds up to eight stages applied in order;
Table~\ref{tab:operators}
lists the operators. Point operators (grayscale, invert, threshold, tint, light)
are composed into one 256-entry lookup table per plane for each run
of consecutive
point stages, so that a run costs one pass over the frame whatever
its length. The
neighbourhood operators read 3$\times$3 or 5-sample neighbourhoods
and replicate the
border samples, because zero padding would produce a false response
along the frame
edge.

The Gaussian blur repeats a separable binomial kernel,
\begin{equation}\label{eq:blur}
  g(x) = \left\lfloor \frac{p_{x-2} + 4p_{x-1} + 6p_x + 4p_{x+1} +
  p_{x+2} + 8}{16}
  \right\rfloor ,
\end{equation}
horizontally and then vertically, $s$ times for strength $s$;
repeated passes widen
the effective Gaussian as $\sqrt{s}$. Each thread computes its own
band of output rows
in a single sweep, keeping the five horizontally filtered rows it
needs in a rolling
window, and successive passes alternate between the destination and a
scratch plane,
which needs one synchronization of the thread team per pass where
filtering in place
needs two. From $s = 4$ the passes run at half resolution: the plane
is reduced by
2$\times$2 averaging, blurred, and expanded by bilinear interpolation
with weights
$\tfrac34$ and $\tfrac14$.

Sobel \cite{duda1973} computes the gradients $G_x$ and $G_y$ of the
3$\times$3 Sobel
kernels and outputs $\min(|G_x| + |G_y|, 255)$, set to zero below a threshold.
The Laplacian-of-Gaussian operator \cite{marr1980} blurs with
Eq.~\eqref{eq:blur}, takes the 4-neighbour Laplacian in 16-bit signed
arithmetic and
marks a pixel where the response changes sign towards its right or
lower neighbour by
at least a threshold; comparing in two directions only keeps contours
one pixel wide,
and the slope test rejects noise that merely grazes zero. Canny's detector
\cite{canny1986} is implemented in four stages: at least one blur
pass, which keeps
the cost of hysteresis bounded on noisy input; the Sobel gradient,
with its direction
quantized to four sectors by integer comparisons ($|G_y| < 0.4|G_x|$ horizontal,
  $|G_y| > 2.4|G_x|$ vertical, otherwise diagonal by the signs of $G_x$
and $G_y$);
non-maximum suppression, keeping a pixel when it is not below one
neighbour across the
edge and strictly above the other, which thins a symmetric ridge to
one pixel; and
hysteresis. Suppression labels each pixel directly in the output, as
a seed when its
magnitude reaches the high threshold and as a candidate when it
reaches the low one;
each row's seeds are collected as it is labelled; a serial flood fill
from the seeds
promotes connected candidates, and a final parallel pass clears the
candidates it did
not reach. The flood fill gives the exact connected set, which a
parallel sweep would
only approximate.

The edge operators output luma only and set the chroma planes to the
neutral value
128. A pass that computes chroma no later stage will read (a blur,
  tint or contrast
stretch before an edge operator, for example) is skipped; the stages
after the last
one that writes chroma are traced back once per configuration. When a
chain's output
chroma is known to be neutral, the frame is marked and the frame bus
copies only luma
(Section~\ref{sec:bus}).

The neighbourhood operators have NEON implementations, written with
intrinsics that
exist on both ARMv7 and AArch64, and portable scalar versions that the compiler
vectorizes on x86. Loops are parallelized with OpenMP \cite{dagum1998}; every
parallel loop is free of data races, so that it is also correct
without OpenMP. Each
vector or parallel path was required to produce output byte-identical
to the scalar
single-threaded path (Section~\ref{sec:verification}).

\begin{table*}[!t]
  \centering\small
  \caption{Image operators of the chain. All parameters are integers
  from 0 to 255.}\label{tab:operators}
  \setlength{\tabcolsep}{4pt}
  \begin{tabular}{@{}lL{4.2cm}L{6.4cm}l@{}}
    \toprule
    Operator & Parameters (default) & Operation & Planes \\
    \midrule
    grayscale & -- & chroma set to 128 & UV \\
    invert & -- & $Y \leftarrow 255 - Y$ & Y \\
    threshold & level (0) & $Y \leftarrow 255$ if $Y \geq$ level, else 0 & Y \\
    tint & $u$, $v$, strength (0) & chroma moved towards ($u$, $v$)
    by strength/255 & UV \\
    light & level (128) & $Y \leftarrow Y + \text{level} - 128$;
    chroma scaled about 128 by level/128 & YUV \\
    contrast & -- & per-frame min/max stretch of $Y$ to 0--255 & Y \\
    blur & strength $s$ (0, one pass) & Eq.~\eqref{eq:blur}, $s$
    passes; half resolution from $s = 4$ & YUV \\
    sobel & threshold (0) & $\min(|G_x|+|G_y|, 255)$, zero below
    threshold & Y; UV $\leftarrow$ 128 \\
    log & strength (0, one pass), threshold (0) & blur, 4-neighbour
    Laplacian, zero-crossings with slope test & Y; UV $\leftarrow$ 128 \\
    canny & strength (0, one pass), low (40), high (90) & blur,
    gradient, suppression, hysteresis & Y; UV $\leftarrow$ 128 \\
    \bottomrule
  \end{tabular}
  \tablenote{Consecutive point operators (the first five rows) are
    fused into one lookup
    table per plane. The image statistics and the visibility correction of
  Section~\ref{sec:autovis} run before the chain and are not chain stages.}
\end{table*}
